\chapter{Matematiğin Doğası ve Matematiksel Düşünme}
\label{chap:matematigin-dogasi}


\subsection*{Öğrenme Hedefleri}

Bu bölümü tamamladığınızda:

\begin{itemize}\setlength{\itemsep}{2pt}\setlength{\parskip}{0pt}
\item Matematik felsefesinin ne olduğunu ve hangi soruları sorduğunu açıklayabileceksiniz.
\item Matematiksel ontoloji ile matematiksel epistemoloji arasındaki farkı ayırt edebileceksiniz.
\item Matematiğin “gerçek dünyaya şaşırtıcı biçimde iyi uygulanabilirliği” probleminin neden felsefi bir problem olduğunu tartışabileceksiniz.
\item Muhakeme, mantıksal dil, ispat ve matematiksel yapılar arasındaki ilişkiyi kitabın genel planı içinde açıklayabileceksiniz.
\end{itemize}

\section{Matematik Neden Felsefi Bir Soru Doğurur?}

Çoğu insan için matematik, kesin, tartışmasız ve şüpheye yer bırakmayan bir bilgi alanı olarak görülür. “2 + 2 = 4” önermesinin doğruluğundan kimse şüphe etmez; bir üçgenin iç açıları toplamının belirli koşullar altında 180 derece olduğunu ispatlayabiliriz ve bu ispat, bir kere yapıldıktan sonra, kabul edilen aksiyomlar ve çıkarım kuralları içinde kalındığı sürece geçerliliğini korur. Matematiğin bu kesinliği, onu diğer bilgi alanlarından ayıran en çarpıcı özelliklerinden biridir. Ancak bu, matematiksel tanımların, aksiyomların ve uygulamaların felsefi bakımdan tartışmasız olduğu anlamına gelmez (Shapiro 2000; Colyvan 2012). Ancak bu kesinlik, daha derin bir soruşturmaya girildiğinde, aslında hiç de basit olmayan sorular doğurur. Sayılar nerede bulunur? “3” sayısı fiziksel evrende bir yerde mi vardır, yoksa yalnızca insan zihninin bir kurgusu mudur? Bir matematikçi yeni bir teoremi keşfeder mi (sanki önceden var olan bir şeyi bulmuş gibi), yoksa onu icat mı eder (sanki daha önce var olmayan bir şeyi yaratmış gibi)? Sonsuz

kümeler gerçekten var mıdır, yoksa sonsuzluk yalnızca zihnimizin bir kurgusu, bir “as if” (sanki-varmış-gibi) düşünme biçimi midir? İşte matematik felsefesi, tam olarak bu tür soruları sistematik biçimde ele alan bir disiplindir. Matematik felsefesi, matematiği yapmakla değil, matematiğin ne olduğunu anlamakla ilgilenir. Bir matematikçi bir teoremi ispatlamakla uğraşırken, matematik filozofu “ispat nedir, bir şeyin ispatlanmış olması ne anlama gelir, matematiksel nesneler var mıdır, matematik neden fiziksel dünyaya bu kadar başarılı uygulanır” gibi sorularla uğraşır (Benacerraf ve Putnam 1983; Shapiro 2000). Bu iki uğraşın birbirinden bağımsız olmadığını da baştan belirtmek gerekir. Matematik tarihine bakıldığında, en derin felsefi tartışmaların çoğu zaman matematiğin içinden, yani bizzat matematikçilerin karşılaştığı somut güçlüklerden doğduğu görülür. Bir matematikçi yeni bir sayı türüyle (örneğin irrasyonel ya da sanal sayılarla) karşılaştığında, önce bu yeni nesnenin nasıl işlem göreceğini bulmaya çalışır; ama er ya da geç “bu nesne gerçekten var mı, yoksa yalnızca hesaplamayı kolaylaştıran kâğıt üzerinde bir kurgu mu?” sorusu da gündeme gelir. Bu ders boyunca göreceğimiz gibi, matematik felsefesi soyut bir “dışarıdan bakış” değil, matematiğin kendi tarihsel gelişiminin doğal bir uzantısıdır.

\begin{tcolorbox}[bookbox,colback=cEx]\textbf{\textsf{\textcolor{cAcc}{Örnek.}}}\quad \textbf{Pisagorcu anlatı ve ortak ölçülemezlik.} Geleneksel anlatıda Pisagorcuların kozmik düzeni sayılar ve basit oranlarla açıkladıkları, karenin kenarı ile köşegeninin ortak ölçülemezliğinin ise bu görüşle kuramsal bir gerilim yarattığı söylenir. Bugünkü dille kenarı 1 olan karenin köşegeni $\sqrt2$’dir ve hiçbir tam sayı oranıyla ifade edilemez. Ancak bu sonucu ilk kimin, hangi ispatla elde ettiği ve erken Pisagorcu çevrede belgelenmiş ani bir “felsefi şok” yaşanıp yaşanmadığı bilinmemektedir. Bu nedenle örnek, tek bir tarihsel olayın kaydı olarak değil, matematiksel bir sonucun daha geniş bir kavramsal çerçeveyle nasıl gerilim kurabileceğini gösteren geleneksel bir anlatı olarak okunmalıdır (Fowler 1999; Huffman 2019).\end{tcolorbox}

\section{Matematiksel Ontoloji ve Epistemoloji}

\begin{tcolorbox}[bookbox,colback=cDef]
\textbf{\textsf{\textcolor{cAcc}{Tanım.}}}\quad
\emph{Matematiksel ontoloji}, matematiksel nesnelerin ne olduğu ve nasıl var olduğu sorusunu; \emph{matematiksel epistemoloji} ise bu nesne ve doğrular hakkında nasıl bilgi edindiğimiz sorusunu inceler.
\end{tcolorbox}

Matematik felsefesindeki sorular genellikle iki büyük başlık altında toplanır. Ontoloji (varlık bilgisi), var olan şeylerin ne olduğunu ve nasıl var olduğunu inceleyen felsefe dalıdır. Matematiksel ontoloji ise şu soruyu sorar: Matematiksel nesneler (sayılar, kümeler, fonksiyonlar, geometrik şekiller) var mıdır ve varsalar nasıl bir varlığa sahiptirler? Bu soruya verilen cevaplar büyük ölçüde farklılaşır. Platonist bir görüşe göre sayılar, zaman ve mekândan bağımsız, insan zihninden önce var

olan soyut nesnelerdir — tıpkı bir gezegenin varlığını keşfetmemiz gibi, matematikçiler de sayıların özelliklerini keşfeder. Çağdaş matematiksel platonizm genellikle varlık, soyutluk ve insan düşüncesinden bağımsızlık tezleriyle tanımlanır; tarihsel Platon’un görüşüyle birebir özdeş değildir. Buna karşılık nominalist yaklaşımlar, soyut matematiksel nesnelere ontolojik bağlılığı reddeder veya matematiksel söylemin böyle bir bağlılık gerektirmediğini göstermeye çalışır. Nominalizmi “sembollerle oynanan kurallı oyun” görüşüyle özdeşleştirmek doğru değildir; bu ifade belirli biçimcilik yorumlarına daha yakındır (Bueno 2020; Field 1980). Bu iki uç arasında pek çok ara pozisyon da mevcuttur. Epistemoloji (bilgi bilgisi) ise bilginin doğasını, kaynağını ve sınırlarını inceler. Matematiksel epistemoloji şu soruyu sorar: Matematiksel bilgiye nasıl ulaşırız ve bu bilgiye neden bu kadar güvenebiliriz? Deneysel bilimlerde bilgi, gözlem ve deneyle test edilir; yanlışlanabilir ve geçicidir. Matematiksel bilgi ise akıl yürütme ve ispat yoluyla elde edilir ve bir kere ispatlandıktan sonra —eğer ispat doğruysa— bir daha geçersiz kılınamaz gibi görünür. Bu, matematiksel bilginin neden bu denli “kesin” göründüğünü açıklamaya çalışan felsefi bir problemdir. Bu iki soru kümesi (ontoloji ve epistemoloji) aslında birbirine sıkı sıkıya bağlıdır: matematiksel nesnelerin ne tür bir varlığa sahip olduğuna dair görüşümüz, bu nesneler hakkında nasıl bilgi edinebileceğimize dair görüşümüzü de şekillendirir. Örneğin, sayıların soyut ve zihin-bağımsız nesneler olduğunu savunan bir platonist, bu nesnelere nasıl “erişebildiğimizi” açıklamakta zorlanır — çünkü bu nesneler fiziksel dünyada değildir, onlarla duyularımız aracılığıyla etkileşime giremeyiz. Bu zorluk, felsefede epistemolojik erişim problemi olarak anılır ve Paul Benacerraf’ın 1973 tarihli “Mathematical Truth” makalesinde, matematik için standart bir doğruluk semantiği ile dönemin nedensel bilgi anlayışını birlikte karşılamanın güçlüğü olarak keskinleştirilmiştir (Benacerraf 1973, ss. 661–668). Soyut ve nedensel bakımdan etkisiz nesneler hakkındaki inançlarımızın doğruluğunu nasıl açıkladığımız, platonist yaklaşımlar için önemli fakat farklı cevaplara açık bir sorundur (Linnebo 2023). Benacerraf’ın 1965 tarihli “What Numbers Could Not Be” makalesi ise başka bir problemi ele alır: doğal sayılar farklı fakat eşit ölçüde uygun küme kuramsal nesnelerle temsil edilebiliyorsa, örneğin “3”ü bu temsillerden yalnızca biriyle özdeşleştirmek için matematiksel bir neden bulmak güçleşir. Bu özdeşleştirme problemi, 1973’teki epistemolojik problemle karıştırılmamalıdır (Benacerraf 1965, ss. 47–62). Bu iki tartışma, çağdaş matematik felsefesinin gündeminin merkezinde yer alır. Ontoloji ve epistemoloji arasındaki ilişkiyi somutlaştırmanın bir yolu,

iki farklı sorunun aynı olguyu farklı açılardan sorduğunu görmektir: “Asal sayılar sonsuzdur” önermesi ne hakkında ve neye dayanarak doğrudur (bu, ontolojik ve semantik boyutları olan bir sorudur) ve biz bu önermenin doğru olduğunu nereden biliyoruz (bu bir epistemolojik sorudur — bilgiye ulaşma yöntemimiz). İlginç olan, matematikte bu iki sorunun ispat pratiğinde yakından ilişkili olmasıdır: Öklid’in asal sayıların sonsuzluğuna dair ispatı, kabul edilen mantık ve başlangıç ilkeleri altında önermeyi gerekçelendirir ve bu doğruluğu bilme yöntemimizin merkezinde yer alır. Bu yakınlık, matematiği diğer bilgi alanlarından ayıran bir başka özelliktir; örneğin tarihte “Sezar Rubicon’u geçti” önermesinin doğruluğu ile bu önermeyi bilme yöntemimiz (tarihsel kayıtlar, arkeolojik bulgular) birbirinden çok daha ayrıktır.

\begin{figure}[htbp]
\centering
\begin{tikzpicture}[node distance=6mm]
  \node[schemeaccent,text width=8.9cm] (claim) {\textbf{Matematiksel önerme:} “Asal sayılar sonsuzdur.”};
  \node[schemenode,text width=4.05cm,below=of claim,xshift=-2.2cm] (ontology) {\textbf{Ontoloji}\\Ne vardır? Önerme ne hakkındadır ve doğruluğu neye dayanır?};
  \node[schemenode,text width=4.05cm,below=of claim,xshift=2.2cm] (epistemology) {\textbf{Epistemoloji}\\Bu sonuca nasıl ulaşır ve bilgimize neden güveniriz?};
  \node[schemeaccent,text width=8.9cm,below=7mm of $(ontology.south)!0.5!(epistemology.south)$] (proof) {\textbf{İspat: iki soru arasındaki köprü}\\Öklid’in ispatı, önermeyi gerekçelendirirken onu bilme yöntemimizi de oluşturur.};
  \draw[schemearrow] (claim) -- (ontology);
  \draw[schemearrow] (claim) -- (epistemology);
  \draw[schemearrow] (ontology) -- (proof);
  \draw[schemearrow] (epistemology) -- (proof);
\end{tikzpicture}
\caption{Aynı matematiksel önermeye ontolojik ve epistemolojik açılardan yaklaşmak.}
\end{figure}

\begin{tcolorbox}[bookbox,colback=cEx]\textbf{\textsf{\textcolor{cAcc}{Örnek.}}}\quad \textbf{Sayı “3” nerede?} Masanın üzerinde üç elma olduğunu söylediğimizde, elmaları görebilir, dokunabiliriz. Ama “3” sayısının kendisini nerede görebiliriz? Üç elmayı alıp yesek bile “3” sayısı yok olmaz; başka üç nesneye de uygulanabilir. Bu basit gözlem, sayıların fiziksel nesnelerden farklı bir varlık türüne sahip olabileceğini düşündürür ve matematiksel ontolojinin çıkış noktalarından biridir.\end{tcolorbox}

\begin{tcolorbox}[bookbox,colback=cEx]\textbf{\textsf{\textcolor{cAcc}{Bilgi Kutusu.}}}\quad Gerçekçilik (Realizm) ile Karşıt Görüşler (Anti-Realizm) Matematik felsefesindeki görüşleri tek bir gerçekçilik–anti-gerçekçilik çizgisine yerleştirmek yararlı fakat kaba bir başlangıçtır. Nesnelerin varlığına ilişkin nesne gerçekçiliği, matematiksel cümlelerin nesnel doğruluk değerlerine ilişkin doğruluk-değeri gerçekçiliği ve matematiksel bilginin zihinden bağımsızlığı birbirinden ayrılmalıdır. Platonizm, nominalizm, kurgusalcılık ve sezgicilik de kendi içlerinde farklı biçimlere sahiptir; örneğin sezgicilik matematiksel nesneleri tümüyle reddetmekten çok onları zihinsel inşaya bağlayabilir. Bu nedenle felsefi ekollere ilişkin sınıflandırmalar kesin kutular değil, karşılaştırma araçları olarak okunmalıdır (Linnebo 2023; Bueno 2020; Shapiro 2000).\end{tcolorbox}

\section{Matematiğin “Mantıksız Derecede” Etkili Olması Problemi}

Matematik felsefesinin bir diğer temel sorusu, fizikçi Eugene Wigner’in 1960 yılında yazdığı ünlü bir makalede “matematiğin doğa bilimlerindeki mantıksız derecede etkinliği” (the unreasonable effectiveness of mathematics) olarak adlandırdığı olgudur. Soyut ve saf akıl yürütmeyle, çoğu zaman herhangi bir pratik uygulama düşünülmeden geliştirilen matematiksel yapılar (örneğin karmaşık sayılar, Öklid dışı geometriler, grup teorisi), yıllar hatta yüzyıllar sonra fizikte, mühendislikte veya başka bilimlerde önemli araçlar hâline gelmiştir. Örneğin 19. yüzyılda genel görelilik amacıyla geliştirilmemiş olan Öklid dışı geometri ve diferansiyel geometri çalışmaları, 20. yüzyılın başında Einstein’ın genel görelilik kuramının matematiksel dilinin parçası hâline gelmiştir (Einstein 1916). Bu tür örnekler yalnızca fizikle sınırlı değildir. 19. yüzyılda tamamen soyut bir sayı teorisi konusu olan asal sayıların dağılımı ve büyük sayıların çarpanlara ayrılmasının zorluğu, 20. yüzyılda açık anahtarlı kriptografinin önemli örneklerinden RSA sisteminde doğrudan kullanılmıştır (Rivest ve diğ. 1978). Benzer biçimde, 19. yüzyılda gelişen grup teorisi de 20. yüzyıl fiziğinde simetrileri betimleyen önemli araçlardan biri hâline gelmiştir. Bu örneklerin ortak noktası, ilgili matematiksel yapıların en azından önemli bölümlerinin sonraki uygulamalarından bağımsız araştırma amaçlarıyla geliştirilmiş olmasıdır; bütün tarihsel aktörlerin hiçbir uygulama düşünmediği gibi daha güçlü bir iddia savunulmamalıdır (Wigner 1960). Bu durum şu soruyu doğurur: Eğer matematik yalnızca insan zihninin keyfi bir icadıysa, fiziksel evrenin yapısını tahmin etmede neden bu kadar başarılı olsun? Bu soru, matematiğin “keşif mi icat mı” olduğu tartışmasını doğrudan besler ve kitabın ilerleyen bölümlerinde

defalarca karşımıza çıkacaktır. Wigner’in sorusuna verilen yanıtlar tek tip değildir. Seçilim etkisi, matematiğin fiziksel sorunlara göre geliştirilmesi, bilim insanlarının uygun matematiksel yapıları etkin biçimde seçmesi ve matematiksel modellerin yalnızca sınırlı yönlerden dünyaya uyması gibi açıklamalar, “etkinliğin” ne kadar şaşırtıcı olduğu konusunda daha ihtiyatlı olmayı gerektirir. Buna karşılık, önceden geliştirilmiş yapıların yeni bilimlerde ayrıntılı ve beklenmedik biçimde işe yaraması sorunun tamamen ortadan kalkmadığını düşündürür. Dolayısıyla Wigner’in makalesi çözülmüş bir teorem değil, farklı felsefi açıklamalara açık bir problem gündemi olarak okunmalıdır (Wigner 1960; Colyvan 2012).

\section{Felsefede Matematiğin Yeri: Kısa Bir Tarihçe}

Matematik ile felsefe arasındaki ilişki, felsefenin kendisi kadar eskidir. Geç dönem kaynaklarında Platon’un Akademi’sinin girişine “Geometri bilmeyen giremesin” yazdırdığı anlatılır; ancak Platon’un yaşadığı dönemden bu yazıtı doğrulayan bir kanıt bulunmadığından, söz tarihsel bir olgu değil, Platoncu eğitimin ruhunu özetleyen sonraki bir gelenek olarak değerlendirilmelidir (Franklin 2023, ss. 365–366). Platon’un matematiğe verdiği felsefi önem, değişmeyen ve akılla kavranan nesneler ile duyusal dünya arasındaki ayrımla bağlantılıdır. Bununla birlikte çağdaş “matematiksel platonizm”, tarihsel Platon’un metinlerinden esinlense de onlarla özdeş bir öğreti değildir (Linnebo 2023). Aristoteles ise matematikçinin duyulur nesneleri hareket ve maddeden “düşüncede ayırarak” nicelik ve şekil bakımından incelediğini savunur. Bu görüş, matematiksel nesneleri basitçe öznel kurgulara indirgemeden, matematiksel soyutlamayı fiziksel nesnelerin özellikleri üzerinden açıklamaya çalışır (Lear 1982, ss. 161–176). Dolayısıyla Aristotelesçi yaklaşım, Platonculuğun yalnızca “zihinsel kurgu” karşıtı olarak özetlenemeyecek daha incelikli bir alternatiftir. 17. yüzyılda Descartes ve Leibniz gibi filozoflar, matematiksel yöntemi kesin bilgi tartışmalarının başlıca modellerinden biri olarak ele aldılar. 18. yüzyılın sonunda Immanuel Kant, Arı Aklın Eleştirisi adlı eserinde matematiksel bilgiyi, ne salt deneyimden (a posteriori) ne de salt tanımlardan (analitik) türeyen, kendine özgü bir bilgi türü —“sentetik a priori”— olarak sınıflandırdı. Kant’a göre “7 + 5 = 12” önermesi, “12” kavramının “7” ve “5” kavramlarının basit bir çözümlemesinden çıkmaz (bu anlamda sentetiktir, yani bilgi verici ve bilgimizi genişleticidir), ama yine de deneyimden bağımsız olarak, kesin ve zorunlu biçimde bilinir (bu anlamda a prioridir) (Kant 1998, B14–B17). Öklid dışı geometrilerin geliştirilmesi, Kant’ın geometri anlayışının

özellikle fiziksel uzay ve zorunluluk bakımından nasıl yorumlanması gerektiği konusunda ciddi bir tartışma doğurmuştur. Bunun Kant’ın bütün sentetik a priori tezini tek başına çürüttüğünü söylemek fazla güçlü olur; saf geometri, insan sezgisinin yapısı ve fiziksel uzayın geometrisi arasındaki ayrımlar dikkate alınmalıdır (Greenberg 2008). 19. yüzyılın sonu ile 20. yüzyılın başı, matematik felsefesi için belki de en verimli dönem olmuştur. Bu dönemde Frege, Russell, Hilbert, Brouwer ve Gödel gibi düşünürler, matematiğin temellerine ilişkin tartışmaları dönüştürmüştür. Bu dönemde ortaya çıkan temel programlar (mantıkçılık, biçimcilik, sezgicilik) farklı tarihlerde, farklı sorunlara yanıt olarak gelişmiş ve kendi içlerinde de çeşitlenmiştir (Benacerraf ve Putnam 1983; Shapiro 2000). Bu kitap ise bu tartışmaların arka planını akılda tutarak matematiksel düşünmenin araçlarına —muhakemeye, mantıksal dile, ispata ve matematiksel yapılara— odaklanacaktır.

\section{Bu Dersin Amacı: Matematik Olmayan Bölümler İçin Neden Önemli?}

Bu noktada haklı bir soru sorulabilir: Matematiksel düşünme ve muhakeme üzerine lisansüstü düzeyde bir ders, bir matematik öğretmeni ya da matematik eğitimi araştırmacısı için ne ifade eder? Bu sorunun üç yanıtı vardır. Birincisi, bu ders bir hesaplama dersi değil, bir düşünme dersidir. Amaç yeni matematiksel teknikler öğretmekten çok, matematiğin arkasındaki mantıksal ve kavramsal yapıyı görünür kılmaktır: bir önermenin ne zaman “ispatlanmış” sayıldığını, bir çıkarımın ne zaman geçerli olduğunu, bir varsayımın ne zaman gizlice devreye girdiğini fark edebilmek. İkincisi, öğretmenin kendi muhakemesi ile öğrencinin muhakemesini çözümleyebilme becerisi birbirinden ayrılamaz. Bir öğrencinin “beş örnekte doğru, öyleyse her zaman doğru” gerekçesini, bir niceleyici sırası hatasını ya da sonucu onaylama yanılgısını teşhis edebilmek, bu yapıları önce kendi düşüncemizde açıkça görebilmeyi gerektirir. Bununla birlikte mantık ve ispat bilgisinin sınıf uygulamasına kendiliğinden aktarılacağını varsaymamak gerekir. Eleştirel düşünme araştırmaları, açık öğretim, tartışma, gerçekçi problemler ve rehberli uygulamanın aktarım için önemli olduğunu göstermektedir (Abrami ve diğ. 2015). Üçüncüsü, kümeler, bağıntılar, kongrüanslar, fonksiyonlar ve kardinalite gibi yapılar, okul matematiğindeki pek çok konunun arkasındaki ortak dili oluşturur. Bu yapılar üzerinde muhakeme etmek, okul konularını birbirine bağlayan kavramsal ağı görünür kılar; kesirlerin denkliği bir denklik bağıntısı, saat aritmetiği bir kongrüans, “her girdiye tek çıktı” kuralı ise fonksiyon kavramıdır.

\begin{tcolorbox}[bookbox,colback=cEx]\textbf{\textsf{\textcolor{cAcc}{Bilgi Kutusu.}}}\quad “Matematik Kaygısı” ve Öğretmenin Rolü Pek çok öğrenci, matematik derslerine geçmiş deneyimlerinden kaynaklanan bir kaygıyla yaklaşabilir. Matematik kaygısı ile başarı arasındaki negatif ilişki iyi belgelenmiştir; fakat ilişki tek nedenli değildir ve her öğrenci için aynı biçimde işlemez (Barroso ve diğ. 2021, ss. 134–168). Bu kitap, öğretmenlere ve öğretmen adaylarına bilinçli olarak farklı bir model sunmayı amaçlar: hızlı işlem ve ezber yerine yavaş, dikkatli ve gerekçeye dayalı düşünme. Kendi ispat denemelerindeki belirsizlik ve takılma anlarını fark eden bir öğretmen, öğrencisinin benzer anlarını da daha iyi tanıyabilir ve onlara daha sabırlı bir yanıt verebilir.\end{tcolorbox}

\section{Dersin Yol Haritası}

Bu kitap, matematiksel düşünmeyi beş aşamada ele alır. İlk aşamada (Bölüm~\ref{chap:matematigin-dogasi}--\ref{chap:muhakeme-ogretimi}) matematiğin doğasını, muhakeme türlerini (tümdengelim, tümevarım, abdüksiyon ve analoji) ve düşünme düzeylerini inceleyecek; muhakemenin nasıl öğretilebileceğini tartışacağız. İkinci aşamada (Bölüm~\ref{chap:onermeler-mantigi}--\ref{chap:niceleyiciler}) önermeler mantığını, argüman geçerliliğini, değişkenleri ve niceleyicileri ortak bir dil içinde kuracağız.

Üçüncü aşamada (Bölüm~\ref{chap:dogrudan-ispat}--\ref{chap:tumevarim}) doğrudan, kontrapozitif, çelişki ve tümevarım yöntemlerini ele alacağız. Dördüncü aşamada (Bölüm~\ref{chap:kumeler}--\ref{chap:kardinalite}) kümelerden kardinaliteye uzanan temel matematiksel yapıları bu yöntemlerle inceleyeceğiz. Son bölümde (Bölüm~\ref{chap:sentez}) kitabın sentezini yapacak ve bu araçların matematik eğitimine yansımalarını değerlendireceğiz. Bu aşamalar arasında sıkı bir mantıksal bağ vardır: muhakeme türleri gerekçelendirmenin biçimlerini ayırt eder; mantık bu biçimleri kesin olarak ifade etmenin dilini verir; ispat yöntemleri iddiaları gerekçelendirir; matematiksel yapılar ise bu yöntemlerin çalışma alanını oluşturur.

\begin{figure}[htbp]
\centering
\begin{tikzpicture}[node distance=4mm]
  \node[schemeaccent,text width=9.4cm] (a) {\textbf{1. Matematiğin doğası ve muhakeme}\\Muhakeme türleri, düşünme düzeyleri ve muhakeme öğretimi};
  \node[schemenode,text width=9.4cm,below=of a] (b) {\textbf{2. Mantıksal dil}\\Önermeler, geçerlilik, değişkenler ve niceleyiciler};
  \node[schemenode,text width=9.4cm,below=of b] (c) {\textbf{3. İspat yöntemleri}\\Doğrudan, kontrapozitif, çelişki ve tümevarım};
  \node[schemenode,text width=9.4cm,below=of c] (d) {\textbf{4. Matematiksel yapılar}\\Kümeler, bağıntılar, kongrüanslar, fonksiyonlar, kardinalite};
  \node[schemeaccent,text width=9.4cm,below=of d] (e) {\textbf{5. Sentez}\\Matematik eğitimine yansımalar ve ders tasarımı};
  \draw[schemearrow] (a)--(b); \draw[schemearrow] (b)--(c);
  \draw[schemearrow] (c)--(d); \draw[schemearrow] (d)--(e);
\end{tikzpicture}
\caption{Kitabın izlediği beş aşamalı kavramsal yol haritası.}
\end{figure}

\section{Yüksek Lisans Derinleştirmesi: Benacerraf Sorunları ve Kuram Seçimi}

Matematik felsefesindeki görüşleri yalnızca “sayılar var mı?” sorusuna verdikleri cevapla sınıflandırmak yetersizdir. Bir konumun, matematiksel nesnelerin kimliği, matematiksel cümlelerin doğruluk koşulları ve bu doğrulara erişimimiz hakkında birlikte işleyen bir açıklama sunması gerekir.

\textbf{Bu kısmın hedefleri.} Bu kısmı tamamladığınızda Benacerraf’ın kimlik ve bilgi sorunlarını ayırabilecek, vazgeçilmezlik argümanını öncül–sonuç biçiminde yeniden kurabilecek, rakip matematik ontolojilerini ortak ölçütlerle karşılaştırabilecek ve bir felsefi konumun maliyetlerini açıkça değerlendirebileceksiniz.

\subsection{İki Benacerraf Sorunu}
Benacerraf’ın 1965 tarihli tartışması, doğal sayıların belirli kümelerle özdeşleştirilmesindeki çoğulluğa dayanır. Von Neumann temsilinde

0 = ∅,     1 = \{∅\},      2 = \{∅, \{∅\}\} iken başka bir kümesel temsil aynı aritmetik yapıyı farklı kümelerle kurabilir. Her iki temsil Peano yapısına izomorf olduğu hâlde “2 gerçekte hangi kümedir?” sorusuna farklı cevap verir. Sorun, aritmetiğin yapısal içeriğiyle belirli bir küme kuramı kodlamasını özdeşleştirmemiz için aritmetiğin kendisinin yeterli gerekçe sağlamamasıdır (Benacerraf 1965, ss. 47–73). 1973 tarihli sorun farklıdır. Standart semantik, matematiksel cümlelerin nesnelere gönderimde bulunarak doğru veya yanlış olduğunu söyler; standart nedensel bilgi anlayışı ise bilginin bilen özne ile bilinen nesne arasında uygun bir nedensel bağlantı gerektirdiğini varsayar. Soyut nesneler nedensel ilişkilere girmiyorsa matematiksel doğruluk ile matematiksel bilgi arasında bir gerilim doğar. Bu, sayıların hangi kümeler olduğu sorunu değil, soyut nesneler hakkındaki bilgimizin nasıl mümkün olduğu sorunudur (Benacerraf 1973, ss. 661–680).

\textbf{Yaygın yanılgı — ‘Benacerraf tek bir itiraz geliştirmiştir.’} Oysa 1965’teki kimlik sorunu, sayıların belirli nesnelerle özdeşleştirilmesini; 1973’teki semantik–epistemolojik gerilim ise nesnel matematiksel doğruların nasıl bilinebildiğini hedef alır. Bir görüş ilkini yapısal konumlarla karşılayıp ikincisi karşısında hâlâ açıklama borçlu kalabilir.

\subsection{Vazgeçilmezlik Argümanının Mantıksal İskeleti}
Quine–Putnam çizgisindeki vazgeçilmezlik argümanı şu biçimde yeniden kurulabilir:

1. En iyi doğrulanmış bilimsel kuramlarımızın vazgeçilmez biçimde nicelediği varlıklara ontolojik bağlılık göstermeliyiz.

2. En iyi bilimsel kuramlarımız matematiksel varlıklar ve yapılar üzerinde vazgeçilmez biçimde niceler.

3. Öyleyse matematiksel varlıklara ontolojik bağlılık göstermeliyiz.

Birinci öncül doğrulama bütüncülüğüne ve ontolojik tutumluluğa, ikinci öncül ise matematiğin bilimden gerçekten elenip elenemeyeceğine bağlıdır (Quine 1948, ss. 32–33; Putnam 1971; Colyvan 2001). Field’ın fiziksel kuramları nominalist bir dilde yeniden kurma programı özellikle ikinci öncüle meydan okur: matematik yararlı ve hatta çıkarım bakımından koruyucu olabilir, fakat bu durum matematiksel nesnelerin varlığını zorunlu kılmayabilir (Field 1980). Argümanın sonucu doğrudan klasik Platonizm değildir. Sonuç kabul edilse bile nesnelerin bağımsız tözler mi, yapılardaki konumlar mı veya bilimsel kuramların başka türden ontolojik gereklilikleri mi olduğu ayrıca tartışılmalıdır.

\subsection{Karşılaştırmalı Kuram Matrisi}

\par\smallskip\begin{center}{\renewcommand{\arraystretch}{1.25}\footnotesize\begin{tabularx}{\linewidth}{|>{\raggedright\arraybackslash}p{2.10cm}|Y|Y|Y|}\hline
\rowcolor{cAcc}\textcolor{white}{\bfseries Konum} & \textcolor{white}{\bfseries Nesne anlayışı} & \textcolor{white}{\bfseries Bilgi kaynağı} & \textcolor{white}{\bfseries Başlıca açıklama yükü} \\ \hline
\textbf{Platonizm} & Zihin ve dilden bağımsız soyut nesneler & Akıl yürütme, sezgi veya kuramsal erişim & Nedensel olmayan epistemik erişim \\ \hline
\textbf{Yapısalcılık} & Bir yapıdaki konumlar veya ilişkisel örüntüler & Yapısal betimlemeler ve izomorfizme değişmez çıkarımlar & Yapıların ontolojik statüsü \\ \hline
\textbf{Nominalizm} & Soyut matematiksel nesneler reddedilir & Somut sistemler, dil veya nominalist kuramlar & Bilimsel uygulamalardaki matematiğin elenmesi \\ \hline
\textbf{Kurgusalcılık} & Matematiksel söylem yararlı bir kurgu gibi işler & Kurgunun içindeki çıkarımlar ve uygulama köprüleri & Yanlış varsayılan bir kuramın güvenilir uygulaması \\ \hline
\end{tabularx}}\end{center}\par\smallskip

Bu tablo kesin sınıflar değil, araştırma soruları üretmek için bir karşılaştırma aracıdır. Aynı başlık altında farklı epistemoloji ve semantik önerileri bulunabilir (Shapiro 1997; Colyvan 2012).

\subsection{Seminer ve Araştırma Çalışmaları}
1. Benacerraf’ın 1965 ve 1973 sorunlarını, aynı görüşün birine cevap verip diğerinde neden zorlanabileceğini gösteren bir örnekle ayırın.

2. Von Neumann ve alternatif bir doğal sayı temsilinin ilk dört terimini yazın; aritmetik bakımdan değişmeyen ve değişen özellikleri belirleyin.

3. Vazgeçilmezlik argümanının iki öncülünden hangisinin daha tartışmalı olduğunu savunan 1.000 sözcüklük bir değerlendirme yazın.

4. Genel görelilik veya açık anahtarlı kriptografi örneğinde “uygulamada vazgeçilmezlik” ile “ontolojik bağlılık” arasındaki geçişi çözümleyin.

5. Platonizm, yapısalcılık ve nominalizmden ikisini seçerek doğruluk, bilgi ve uygulama ölçütleri bakımından karşılaştırmalı bir araştırma önerisi hazırlayın.

\section*{Bölüm Özeti}

Matematik felsefesi, matematiğin ne olduğunu, matematiksel nesnelerin nasıl var olduğunu (ontoloji) ve matematiksel bilginin nasıl mümkün olduğunu (epistemoloji) inceleyen bir disiplindir. Matematiğin fiziksel dünyaya şaşırtıcı derecede iyi uygulanabilmesi, “keşif mi icat mı” tartışmasını doğuran temel bir problemdir. Bu kitap, matematiğin doğasına ilişkin bu soruların üzerine muhakeme türlerini, mantıksal dili, ispat yöntemlerini ve temel matematiksel yapıları sırasıyla kuracaktır.

\section*{Gözden Geçirme Soruları}

\par\noindent 1. Matematiksel ontoloji ile matematiksel epistemoloji arasındaki farkı kendi cümlelerinizle açıklayınız.

\par\noindent 2. “Matematikçiler teoremleri keşfeder mi, icat mı ederler?” sorusuna kendi görüşünüzü gerekçelendirerek tartışınız.

\par\noindent 3. Wigner’in “matematiğin mantıksız derecede etkinliği” tezini bir örnekle açıklayınız.

\par\noindent 4. Günlük hayattan, matematiğin beklenmedik biçimde işe yaradığı bir örnek düşününüz ve bunun neden şaşırtıcı olduğunu tartışınız.

\section*{Çoktan Seçmeli Sorular}

\par\smallskip\noindent 1. Matematik felsefesinin temel uğraşı aşağıdakilerden hangisidir?

\par\noindent\hspace*{1.6em}A) Yeni matematiksel teoremler ispatlamak

\par\noindent\hspace*{1.6em}B) Matematiğin ne olduğunu, doğasını ve temellerini sorgulamak

\par\noindent\hspace*{1.6em}C) Öğrencilere hızlı hesaplama teknikleri öğretmek

\par\noindent\hspace*{1.6em}D) Matematik tarihindeki tüm formülleri kronolojik olarak listelemek

\par\smallskip\noindent 2. Aşağıdakilerden hangisi matematiksel ontoloji sorusuna bir örnektir?

\par\noindent\hspace*{1.6em}A) “Bu ispat yöntemi ne kadar sürede tamamlanır?”

\par\noindent\hspace*{1.6em}B) “Sayılar insan zihninden bağımsız olarak var mıdır?”

\par\noindent\hspace*{1.6em}C) “Bu teoremi kim ilk kez ispatlamıştır?”

\par\noindent\hspace*{1.6em}D) “Bu formül hangi mühendislik alanında kullanılır?”

\par\smallskip\noindent 3. Aşağıdakilerden hangisi matematiksel epistemoloji sorusuna bir örnektir?

\par\noindent\hspace*{1.6em}A) “Küme kavramı ne zaman ortaya çıkmıştır?”

\par\noindent\hspace*{1.6em}B) “Matematiksel bilgiye nasıl ulaşırız ve bu bilgiye neden güvenebiliriz?”

\par\noindent\hspace*{1.6em}C) “Bu geometrik şeklin alanı nasıl hesaplanır?”

\par\noindent\hspace*{1.6em}D) “Bu teorem hangi dilde ilk kez yayımlanmıştır?”

\par\smallskip\noindent 4. Platonist görüşe göre matematiksel nesnelerle ilgili aşağıdaki ifadelerden hangisi doğrudur?

\par\noindent\hspace*{1.6em}A) Matematiksel nesneler yalnızca sembollerle oynanan bir oyunun parçalarıdır.

\par\noindent\hspace*{1.6em}B) Matematiksel nesneler insan zihninin keyfi icatlarıdır.

\par\noindent\hspace*{1.6em}C) Matematiksel nesneler zaman ve mekândan bağımsız, insan zihninden önce var olan soyut varlıklardır.

\par\noindent\hspace*{1.6em}D) Matematiksel nesneler yalnızca fiziksel nesnelerin isimlendirilmesinden ibarettir.

\par\smallskip\noindent 5. Nominalist görüşe göre matematiksel nesnelerle ilgili aşağıdakilerden hangisi doğrudur?

\par\noindent\hspace*{1.6em}A) Matematiksel nesneler bağımsız bir varlık alanında gerçekten vardır.

\par\noindent\hspace*{1.6em}B) Soyut matematiksel nesneler yoktur ya da matematiksel söylem bu tür nesnelere ontolojik bağlılık gerektirmeden yorumlanmalıdır.

\par\noindent\hspace*{1.6em}C) Matematiksel nesneler yalnızca fizik biliminin nesneleriyle özdeştir.

\par\noindent\hspace*{1.6em}D) Matematiksel nesneler yalnızca geometrik şekillerden ibarettir.

\par\smallskip\noindent 6. “Sayı 3’ün kendisini nerede görebiliriz?” örneği hangi felsefi tartışmayı somutlaştırmak için kullanılmıştır?

\par\noindent\hspace*{1.6em}A) İspat yöntemlerinin sınıflandırılması

\par\noindent\hspace*{1.6em}B) Matematiksel nesnelerin varlık türünün fiziksel nesnelerden farklı olabileceği

\par\noindent\hspace*{1.6em}C) Aksiyom sistemlerinin tutarlılığı

\par\noindent\hspace*{1.6em}D) Niceleyicilerin doğru kullanımı

\par\smallskip\noindent 7. Paul Benacerraf’ın öne sürdüğü “epistemolojik erişilebilirlik problemi” temelde neyi sorgular?

\par\noindent\hspace*{1.6em}A) Matematikçilerin ne kadar hızlı hesap yapabildiğini

\par\noindent\hspace*{1.6em}B) Fiziksel bir varlık olan insanın, soyut ve zaman-mekân dışı matematiksel nesneler hakkında nasıl bilgi edinebildiğini

\par\noindent\hspace*{1.6em}C) Matematik derslerinin müfredatta hangi sırayla okutulması gerektiğini

\par\noindent\hspace*{1.6em}D) Öklid dışı geometrilerin fizikte nasıl kullanıldığını

\par\smallskip\noindent 8. Eugene Wigner’in 1960 tarihli makalesinde ortaya attığı tez, aşağıdakilerden hangisidir?

\par\noindent\hspace*{1.6em}A) Matematik, doğa bilimlerine uygulanamayacak kadar soyuttur.

\par\noindent\hspace*{1.6em}B) Matematiğin doğa bilimlerindeki etkinliği “mantıksız derecede” yani beklenenin çok ötesinde başarılıdır.

\par\noindent\hspace*{1.6em}C) Fizik, matematikten bağımsız olarak gelişebilir.

\par\noindent\hspace*{1.6em}D) Matematiksel formüller yalnızca mühendislikte kullanılabilir.

\par\smallskip\noindent 9. Aşağıdaki örneklerden hangisi, Wigner’in tezini destekleyen tarihsel bir örnek olarak metinde verilmiştir?

\par\noindent\hspace*{1.6em}A) Öklid dışı geometrinin ve diferansiyel geometrinin, genel görelilik amacıyla geliştirilmemişken daha sonra Einstein’ın kuramında kullanılması

\par\noindent\hspace*{1.6em}B) Pisagor teoreminin ilkokullarda öğretilmesi

\par\noindent\hspace*{1.6em}C) Kalkülüsün mühendislik fakültelerinde zorunlu ders olması

\par\noindent\hspace*{1.6em}D) İstatistiğin anket çalışmalarında kullanılması

\par\smallskip\noindent 10. Metne göre, sayı teorisinin daha sonra hangi pratik alanda önemli bir araç hâline geldiği belirtilmiştir?

\par\noindent\hspace*{1.6em}A) Tarım mühendisliği

\par\noindent\hspace*{1.6em}B) Modern kriptografi (örneğin RSA şifreleme sistemi)

\par\noindent\hspace*{1.6em}C) Mimari tasarım

\par\noindent\hspace*{1.6em}D) Tıbbi görüntüleme

\par\smallskip\noindent 11. “Matematiğin etkililiği bir mucize değil, bir seçilim etkisidir” biçimindeki itiraz neyi savunur?

\par\noindent\hspace*{1.6em}A) Matematiğin fiziğe uygulanabilirliği yalnızca işe yarayan yapıların hatırlanıp diğerlerinin unutulmasından kaynaklanır.

\par\noindent\hspace*{1.6em}B) Matematik hiçbir zaman fiziğe uygulanamaz.

\par\noindent\hspace*{1.6em}C) Fizikçiler matematiği kullanmayı bırakmalıdır.

\par\noindent\hspace*{1.6em}D) Matematiksel yapılar rastgele seçilerek fiziğe uygulanır.

\par\smallskip\noindent 12. Antik Yunan filozofu Platon’a göre matematiksel nesnelerin konumu nasıldır?

\par\noindent\hspace*{1.6em}A) Fiziksel nesnelerden soyutlama yoluyla zihinde oluşurlar.

\par\noindent\hspace*{1.6em}B) Duyusal dünyanın ötesindeki, mükemmel ve değişmez “İdealar” alanında bulunurlar.

\par\noindent\hspace*{1.6em}C) Yalnızca dilsel semboller olarak vardırlar.

\par\noindent\hspace*{1.6em}D) İnsan deneyiminin bir ürünüdürler ve zamanla değişirler.

\par\smallskip\noindent 13. Aristoteles’in matematiksel nesnelere yaklaşımı, Platon’unkinden temelde hangi noktada ayrılır?

\par\noindent\hspace*{1.6em}A) Aristoteles matematiksel nesnelerin fiziksel nesnelerden soyutlama yoluyla zihinde oluştuğunu, ayrı bir “İdealar dünyasına” gerek olmadığını savunmuştur.

\par\noindent\hspace*{1.6em}B) Aristoteles matematiğin tamamen anlamsız olduğunu savunmuştur.

\par\noindent\hspace*{1.6em}C) Aristoteles matematiksel nesnelerin yalnızca Tanrı’nın zihninde var olduğunu savunmuştur.

\par\noindent\hspace*{1.6em}D) Aristoteles ile Platon’un görüşleri arasında hiçbir fark yoktur.

\par\smallskip\noindent 14. Immanuel Kant’a göre “7 + 5 = 12” gibi matematiksel önermeler hangi bilgi türüne örnektir?

\par\noindent\hspace*{1.6em}A) Analitik a posteriori

\par\noindent\hspace*{1.6em}B) Sentetik a priori

\par\noindent\hspace*{1.6em}C) Analitik a priori

\par\noindent\hspace*{1.6em}D) Sentetik a posteriori

\par\smallskip\noindent 15. Kant’ın geometri anlayışının nasıl yorumlanması gerektiği konusunda hangi gelişme önemli bir tartışma doğurmuştur?

\par\noindent\hspace*{1.6em}A) Kalkülüsün icadıyla

\par\noindent\hspace*{1.6em}B) Öklid dışı geometrilerin keşfiyle

\par\noindent\hspace*{1.6em}C) Bilgisayarların icadıyla

\par\noindent\hspace*{1.6em}D) İstatistik biliminin gelişmesiyle

\par\smallskip\noindent 16. Metne göre matematiksel düşünme ve muhakeme üzerine bir ders, öğrenciye öncelikle ne kazandırmayı amaçlar?

\par\noindent\hspace*{1.6em}A) Öğrencilere ileri düzey hesaplama teknikleri öğretmek için

\par\noindent\hspace*{1.6em}B) Titiz akıl yürütme, kanıt değerlendirme ve varsayımları sorgulama becerilerini geliştirmek için

\par\noindent\hspace*{1.6em}C) Okul matematiğindeki formülleri ezberletmek için

\par\noindent\hspace*{1.6em}D) Sınavlarda hızlı işlem yapmayı öğretmek için

\par\smallskip\noindent 17. Metinde “korelasyon nedensellik değildir” örneği hangi amaçla verilmiştir?

\par\noindent\hspace*{1.6em}A) İstatistik formüllerinin nasıl hesaplandığını göstermek için

\par\noindent\hspace*{1.6em}B) Matematiksel/mantıksal düşünmenin günlük hayattaki akıl yürütme hatalarını fark etmeye nasıl yardımcı olduğunu göstermek için

\par\noindent\hspace*{1.6em}C) Korelasyon katsayısının hesaplanma yöntemini açıklamak için

\par\noindent\hspace*{1.6em}D) Gazetecilik mesleğinin tarihini anlatmak için

\par\smallskip\noindent 18. Kitabın dördüncü aşamasında hangi konular ele alınacaktır?

\par\noindent\hspace*{1.6em}A) Değişken ve niceleyici kavramları

\par\noindent\hspace*{1.6em}B) Muhakeme türleri ve düşünme düzeyleri

\par\noindent\hspace*{1.6em}C) Kümeler, bağıntılar, kongrüanslar, fonksiyonlar ve kardinalite

\par\noindent\hspace*{1.6em}D) Matematik felsefesindeki ekoller

\par\smallskip\noindent 19. Metne göre kitabın bölümleri arasındaki mantıksal bağ nasıl özetlenmiştir?

\par\noindent\hspace*{1.6em}A) Bölümler birbirinden tamamen bağımsız, ayrı konulardır.

\par\noindent\hspace*{1.6em}B) Muhakeme türleri gerekçelendirme biçimlerini ayırt eder, mantık bunları kesin biçimde ifade etmenin dilini verir, ispat yöntemleri iddiaları gerekçelendirir, matematiksel yapılar bu yöntemlerin çalışma alanını oluşturur.

\par\noindent\hspace*{1.6em}C) Bölümler yalnızca tarihsel kronolojiye göre sıralanmıştır ve aralarında kavramsal bir bağ yoktur.

\par\noindent\hspace*{1.6em}D) Her bölüm, bir önceki bölümü tamamen çürütmektedir.

\par\smallskip\noindent 20. Metinde “matematik kaygısı” ile ilgili bilgi kutusunda kitabın önerdiği yaklaşım nasıl tanımlanmıştır?

\par\noindent\hspace*{1.6em}A) Hızlı işlem yapmayı ve formül ezberlemeyi ön planda tutan bir yaklaşım

\par\noindent\hspace*{1.6em}B) Hızlı işlem ve ezber yerine yavaş, dikkatli ve gerekçeye dayalı düşünmeyi ön planda tutan bir yaklaşım

\par\noindent\hspace*{1.6em}C) Yalnızca matematik bölümü öğrencilerine uygun, ileri düzey bir yaklaşım

\par\noindent\hspace*{1.6em}D) Sınav odaklı, ezbere dayalı bir yaklaşım

\clearpage

\clearpage